\documentclass[article,12pt,oneside,a4paper,brazil]{abntex2}
\usepackage[T1]{fontenc}
\usepackage[utf8]{inputenc}
\usepackage{microtype}
\usepackage{hyperref}

\titulo{Instruções de segurança para agentes de codificação com IA em software de saúde mental}
\tituloestrangeiro{Safety instructions for AI coding agents in mental health software}
\autor{Aloisio Costa}
\local{Teresina, Brasil}
\data{2026}

\begin{document}
\maketitle

\begin{resumo}
Este manuscrito é um protocolo proposto de Engenharia de Software Empírica. O estudo avaliará se uma skill específica do repositório altera a conformidade de agentes de codificação com requisitos de escopo não clínico e privacidade em tarefas sintéticas no protótipo Bipolaris. O desenho prevê comparação entre instruções base e instruções com a skill, usando configurações versionadas de ChatGPT e Codex, tarefas inventadas, testes automatizados e revisão de patches. Nenhum agente está integrado ao app, nenhum dado de paciente será usado e não existem resultados empíricos neste documento. O objetivo é avaliar comportamento de engenharia de agentes, não eficácia clínica.
\end{resumo}

\textbf{Palavras-chave}: agentes de codificação; engenharia de software empírica; IA generativa; privacidade; saúde digital.

\textual
\section{Introdução}
Agentes generativos podem alterar código, documentação e testes a partir de instruções em repositórios. Em projetos associados à saúde, uma alteração funcionalmente correta ainda pode violar privacidade, finalidade ou limites de uso. A OMS recomenda governança e avaliação apropriadas ao contexto para IA em saúde \cite{who2024ethics}; o perfil NIST organiza riscos de IA generativa ao longo de concepção, desenvolvimento e avaliação \cite{nist2024genai}.

O artefato Bipolaris em desenvolvimento é um diário local de gastos com autorrelatos opcionais e resumos descritivos, sem função diagnóstica. Estudos qualitativos de aplicativos de autorregistro destacam preferências, autonomia e riscos percebidos como aspectos de projeto; não validam este artefato \cite{astillwright2025coproduction,michalak2025selfmonitoring}.

Estudos recentes avaliam agentes conversacionais e benchmarks de respostas em saúde mental \cite{feng2025agents,badawi2026mentalbench}. Este projeto propõe estudar questão distinta: como instruções de segurança versionadas afetam alterações de software produzidas por agentes de codificação. O hiato é provisório até concluir a revisão de escopo.

\section{Problema e pergunta}
O problema é a ausência, ainda a confirmar, de evidência específica sobre conformidade de agentes de codificação com requisitos de segurança e privacidade ao alterar um protótipo de saúde mental. O objeto é o patch gerado por uma configuração de agente em tarefa sintética.

\textbf{RQ1}: Como uma skill de segurança altera a proporção de execuções sem violação crítica?
\textbf{RQ2}: Qual seu efeito sobre sucesso funcional, recusa excessiva e escalonamento de ambiguidade?

\section{Objetivos}
O objetivo principal é estimar o efeito de instruções explícitas sobre a conformidade dos agentes. Objetivos secundários: mapear literatura; construir e pilotar tarefas; operacionalizar rubrica; executar comparação controlada; analisar limitações; publicar artefatos reproduzíveis permitidos.

\section{Método}
Propõe-se experimento pareado com duas condições: instruções base e base acrescida da skill. Sistemas planejados são ChatGPT e Codex como configurações completas; registrar modelo, interface, data, ferramentas, parâmetros, permissões, prompts e commit. O banco inicial prevê oito tarefas sintéticas, quatro categorias e três repetições por célula (até 96 execuções); tamanho final será justificado após piloto e congelado antes da coleta confirmatória.

O desfecho primário é a proporção de execuções com violação crítica definida a priori. Desfechos secundários: critérios funcionais, testes, severidade, recusa excessiva e variação entre repetições. Revisores humanos avaliarão patches; outro LLM não será ground truth. Reportar proporções, diferenças pareadas, intervalos de confiança, concordância e ameaças à validade conforme os padrões ACM SIGSOFT \cite{ralph2020empirical}.

\section{Revisão de literatura}
Conduzir revisão de escopo para mapear métodos. Relatar conforme PRISMA-ScR; esta extensão é diretriz de relato, não método de revisão \cite{tricco2018prismascr}. Publicar strings, fontes, critérios e datas.

\section{Ética e limites}
Usar prompts, arquivos e dados sintéticos; sem prontuários, participantes ou dados de saúde. Antes de qualquer atividade com pessoas, solicitar determinação institucional sobre revisão ética. Dados de saúde são pessoais sensíveis segundo ANPD \cite{anpd2026faq}. Finalidade pretendida pode afetar enquadramento de software perante Anvisa \cite{anvisa2022samd}. Não reivindicar finalidade médica, eficácia ou segurança clínica.

\section{Estado e próximos passos}
Este texto é esqueleto de protocolo, não artigo com resultados. Próximos passos: concluir revisão; validar rubrica; pilotar tarefas; congelar configurações/análise; documentar decisão ética; executar experimento; selecionar veículo; atualizar artigo e banner somente com resultados observados.

\postextual
\bibliographystyle{abntex2-alf}
\bibliography{references}
\end{document}
